% GENERATED FILE. Edit canonical lessons, then run npm run generate:books.
% canonical-source-hash: fnv1a64:5bbd125e6e78d36b
% canonical-lessons: TA-C473-vacikaram, TA-C473-kampiram, TA-C473-taralam, TA-C473-tannampikkai, TA-C473-meccu

\chapter{Charm and Confidence}
\label{ch:ta-charm-and-confidence}
\hlchaptermodality{473}

\begin{chapteropening}
\textbf{By the end of this chapter:} \emph{I can talk about charm, dignity, generosity and self-confidence.}

Complete the last lesson of chapter 473: I can talk about charm, dignity, generosity and self-confidence.
\end{chapteropening}

\section[\texorpdfstring{\ta{வசீகரம்} (vacīkaram)}{வசீகரம் (vacīkaram)}]{\ta{வசீகரம்} (vacīkaram) — charm; attractiveness}

\label{lesson:TA-C473-vacikaram}



\begin{quote}
Before the new one: say the Tamil for a miser, then the Tamil for a spendthrift.
\end{quote}

\subsection*{You'll want to know: \ta{வசீகரம்}}
\textbf{\ta{வசீகரம்}} — \emph{vacīkaram} — \textquotedblleft{}charm; attractiveness\textquotedblright{}.

\begin{quote}
\ta{அவருடைய} \ta{புன்னகையில்} \ta{ஒரு} \ta{வசீகரம்} \ta{இருந்தது}.
\end{quote}

\begin{quote}
\emph{avaruṭaiya puṉṉakaiyil oru vacīkaram iruntatu.}
\end{quote}

\begin{quote}
There was a charm in his smile.
\end{quote}

\ta{வசீகரமான}, charming.

\begin{cousinweb}
\textbf{\ta{வசீகரம்}} (\emph{vacīkaram}) is borrowed from Sanskrit \emph{vaśīkaraṇa}, \textquotedblleft{}bringing under control; bewitching\textquotedblright{} (Monier-Williams).
\end{cousinweb}

\begin{grammarlens}[title={What you've built}]
One more word for describing what you lived and felt.
\end{grammarlens}

\subsection*{Your turn}
\begin{itemize}
\raggedright
  \item \emph{Say it:} \emph{vacīkaram}
  \item \emph{Say it:} \emph{vacīkaram}, once more
  \item \emph{Say it:} \emph{avaruṭaiya puṉṉakaiyil oru vacīkaram iruntatu.} — There was a charm in his smile.
  \item \emph{Recall:} say the Tamil for a miser, then the Tamil for a spendthrift, then say \emph{vacīkaram} again
\end{itemize}

\begin{tcolorbox}[breakable,colback=teal!4,colframe=teal!35!black,title={Before you move on}]
What does \ta{வசீகரம்} mean? (\textquotedblleft{}Charm; attractiveness\textquotedblright{}.) Say it once more.
\end{tcolorbox}
\section[\texorpdfstring{\ta{கம்பீரம்} (kampīram)}{கம்பீரம் (kampīram)}]{\ta{கம்பீரம்} (kampīram) — majesty; a dignified bearing}

\label{lesson:TA-C473-kampiram}



\begin{quote}
Before the new one: say the Tamil for a spendthrift, then the Tamil for charm.
\end{quote}

\subsection*{You'll want to know: \ta{கம்பீரம்}}
\textbf{\ta{கம்பீரம்}} — \emph{kampīram} — \textquotedblleft{}majesty; a dignified bearing\textquotedblright{}.

\begin{quote}
\ta{யானை} \ta{கம்பீரமாக} \ta{நடந்தது}.
\end{quote}

\begin{quote}
\emph{yāṉai kampīramāka naṭantatu.}
\end{quote}

\begin{quote}
The elephant walked majestically.
\end{quote}

\ta{கம்பீரமான} \ta{குரல்}, a commanding voice.

\begin{cousinweb}
\textbf{\ta{கம்பீரம்}} (\emph{kampīram}) is borrowed from Sanskrit \emph{gambhīra}, \textquotedblleft{}deep\textquotedblright{} (CDIAL 4031); Tamil writes the \emph{g} and \emph{bh} with \ta{க} and \ta{ப}.
\end{cousinweb}

\begin{grammarlens}[title={What you've built}]
One more word for describing what you lived and felt.
\end{grammarlens}

\subsection*{Your turn}
\begin{itemize}
\raggedright
  \item \emph{Say it:} \emph{kampīram}
  \item \emph{Say it:} \emph{kampīram}, once more
  \item \emph{Say it:} \emph{yāṉai kampīramāka naṭantatu.} — The elephant walked majestically.
  \item \emph{Recall:} say the Tamil for a spendthrift, then the Tamil for charm, then say \emph{kampīram} again
\end{itemize}

\begin{tcolorbox}[breakable,colback=teal!4,colframe=teal!35!black,title={Before you move on}]
What does \ta{கம்பீரம்} mean? (\textquotedblleft{}Majesty; a dignified bearing\textquotedblright{}.) Say it once more.
\end{tcolorbox}
\section[\texorpdfstring{\ta{தாராளம்} (tārāḷam)}{தாராளம் (tārāḷam)}]{\ta{தாராளம்} (tārāḷam) — generosity; plenty}

\label{lesson:TA-C473-taralam}



\begin{quote}
Before the new one: say the Tamil for charm, then the Tamil for majesty.
\end{quote}

\subsection*{You'll want to know: \ta{தாராளம்}}
\textbf{\ta{தாராளம்}} — \emph{tārāḷam} — \textquotedblleft{}generosity; plenty\textquotedblright{}.

\begin{quote}
\ta{அவர்} \ta{ஏழைகளுக்குத்} \ta{தாராளமாக} \ta{உதவினார்}.
\end{quote}

\begin{quote}
\emph{avar ēḻaikaḷukkut tārāḷamāka utaviṉār.}
\end{quote}

\begin{quote}
He helped the poor generously.
\end{quote}

\ta{தாராளமாக}, generously; also plenty. There is plenty of room: \ta{தாராளமாக} \ta{இடம்} \ta{இருக்கிறது}.

\begin{cousinweb}
The origin of \textbf{\ta{தாராளம்}} (\emph{tārāḷam}) is not settled; none of our sources lists it.
\end{cousinweb}

\begin{grammarlens}[title={What you've built}]
One more word for describing what you lived and felt.
\end{grammarlens}

\subsection*{Your turn}
\begin{itemize}
\raggedright
  \item \emph{Say it:} \emph{tārāḷam}
  \item \emph{Say it:} \emph{tārāḷam}, once more
  \item \emph{Say it:} \emph{avar ēḻaikaḷukkut tārāḷamāka utaviṉār.} — He helped the poor generously.
  \item \emph{Recall:} say the Tamil for charm, then the Tamil for majesty, then say \emph{tārāḷam} again
\end{itemize}

\begin{tcolorbox}[breakable,colback=teal!4,colframe=teal!35!black,title={Before you move on}]
What does \ta{தாராளம்} mean? (\textquotedblleft{}Generosity; plenty\textquotedblright{}.) Say it once more.
\end{tcolorbox}
\section[\texorpdfstring{\ta{தன்னம்பிக்கை} (taṉṉampikkai)}{தன்னம்பிக்கை (taṉṉampikkai)}]{\ta{தன்னம்பிக்கை} (taṉṉampikkai) — self-confidence}

\label{lesson:TA-C473-tannampikkai}



\begin{quote}
Before the new one: say the Tamil for majesty, then the Tamil for generosity.
\end{quote}

\subsection*{You'll want to know: \ta{தன்னம்பிக்கை}}
\textbf{\ta{தன்னம்பிக்கை}} — \emph{taṉṉampikkai} — \textquotedblleft{}self-confidence\textquotedblright{}.

\begin{quote}
\ta{தேர்வுக்கு} \ta{முன்} \ta{அவளுக்கு} \ta{நிறைய} \ta{தன்னம்பிக்கை} \ta{இருந்தது}.
\end{quote}

\begin{quote}
\emph{tērvukku muṉ avaḷukku niṟaiya taṉṉampikkai iruntatu.}
\end{quote}

\begin{quote}
Before the exam she had plenty of self-confidence.
\end{quote}

\ta{தன்} \textquotedblleft{}oneself\textquotedblright{} + \ta{நம்பிக்கை} \textquotedblleft{}trust\textquotedblright{}: after \ta{ன்} the \ta{ந} becomes \ta{ன}, so the join is written \ta{ன்ன}.

\begin{cousinweb}
\textbf{\ta{தன்னம்பிக்கை}} (\emph{taṉṉampikkai}) is a native compound of \ta{தன்} \emph{taṉ}, \textquotedblleft{}oneself\textquotedblright{} (DEDR 3196), and \ta{நம்பிக்கை}, the noun of \ta{நம்பு} \textquotedblleft{}to trust\textquotedblright{} (DEDR 3600).
\end{cousinweb}

\begin{grammarlens}[title={What you've built}]
One more word for describing what you lived and felt.
\end{grammarlens}

\subsection*{Your turn}
\begin{itemize}
\raggedright
  \item \emph{Say it:} \emph{taṉṉampikkai}
  \item \emph{Say it:} \emph{taṉṉampikkai}, once more
  \item \emph{Say it:} \emph{tērvukku muṉ avaḷukku niṟaiya taṉṉampikkai iruntatu.} — Before the exam she had plenty of self-confidence.
  \item \emph{Recall:} say the Tamil for majesty, then the Tamil for generosity, then say \emph{taṉṉampikkai} again
\end{itemize}

\begin{tcolorbox}[breakable,colback=teal!4,colframe=teal!35!black,title={Before you move on}]
What does \ta{தன்னம்பிக்கை} mean? (\textquotedblleft{}Self-confidence\textquotedblright{}.) Say it once more.
\end{tcolorbox}
\section[\texorpdfstring{\ta{மெச்சு} (meccu)}{மெச்சு (meccu)}]{\ta{மெச்சு} (meccu) — to admire; to commend}

\label{lesson:TA-C473-meccu}



\begin{quote}
Before the new one: say the Tamil for generosity, then the Tamil for self-confidence.
\end{quote}

\subsection*{You'll want to know: \ta{மெச்சு}}
\textbf{\ta{மெச்சு}} — \emph{meccu} — \textquotedblleft{}to admire; to commend\textquotedblright{}.

\begin{quote}
\ta{அவனுடைய} \ta{நேர்மையை} \ta{எல்லோரும்} \ta{மெச்சினார்கள்}.
\end{quote}

\begin{quote}
\emph{avaṉuṭaiya nērmaiyai ellōrum mecciṉārkaḷ.}
\end{quote}

\begin{quote}
Everyone admired his honesty.
\end{quote}

Past (I …): \ta{மெச்சினேன்}, \emph{mecciṉēṉ}.

The quality admired takes -\ta{ஐ}.

\begin{cousinweb}
\textbf{\ta{மெச்சு}} (\emph{meccu}) is native Dravidian (DEDR 4722), where it means to praise and to admire greatly.
\end{cousinweb}

\begin{grammarlens}[title={What you've built}]
One more word for describing what you lived and felt.
\end{grammarlens}

\subsection*{Your turn}
\begin{itemize}
\raggedright
  \item \emph{Say it:} \emph{meccu}
  \item \emph{Say it:} \emph{meccu}, once more
  \item \emph{Say it:} \emph{avaṉuṭaiya nērmaiyai ellōrum mecciṉārkaḷ.} — Everyone admired his honesty.
  \item \emph{Recall:} say the Tamil for generosity, then the Tamil for self-confidence, then say \emph{meccu} again
\end{itemize}

\begin{tcolorbox}[breakable,colback=teal!4,colframe=teal!35!black,title={Before you move on}]
What does \ta{மெச்சு} mean? (\textquotedblleft{}To admire; to commend\textquotedblright{}.) Say it once more.
\end{tcolorbox}
